\documentclass[a4paper,12pt]{article}

% Encodage et langue
\usepackage[utf8]{inputenc}
\usepackage[T1]{fontenc}
\usepackage[french]{babel}

% Mathématiques
\usepackage{amsmath, amssymb, amsfonts}
\usepackage{esint} % double intégrale fermée
\usepackage{physics}
\usepackage{bm} % bold maths

% Graphiques et figures
\usepackage{graphicx}
\usepackage{enumitem} % pour les puces de listes

% Mise en page
\usepackage{geometry}
\geometry{a4paper, margin=1.5cm}
\usepackage{multicol}

% Couleurs et hyperliens
\usepackage{xcolor}
\usepackage{hyperref}
\hypersetup{
    colorlinks=true,
    linkcolor=blue,
    citecolor=blue,
    urlcolor=blue
}

\newcommand{\re}{\textrm}
\newcommand{\Nabla}{\va{\nabla}}
\renewcommand{\curl}{\Nabla \wedge}
\renewcommand{\grad}{\va{\operatorname{grad}}}
\newcommand{\rot}{\va{\operatorname{rot}}}
\renewcommand{\div}{\operatorname{div}}

% Small arrows
\newcommand{\veryshortarrow}[1][4pt]{\mathrel{%
    \hbox{\rule[\dimexpr\fontdimen22\textfont2-.2pt\relax]{#1}{.5pt}}%
    \mkern-4mu\hbox{\usefont{U}{lasy}{m}{n}\symbol{41}}}}
\makeatletter
\setbox0\hbox{$\xdef\scriptratio{\strip@pt\dimexpr
\numexpr(\sf@size*65536)/\f@size sp}$}
\newcommand{\scriptveryshortarrow}[1][4pt]{{%
    \hbox{\rule[\scriptratio\dimexpr\fontdimen22\textfont2-.2pt\relax]
    {\scriptratio\dimexpr#1\relax}{\scriptratio\dimexpr.5pt\relax}}%
    \mkern-4mu\hbox{\let\f@size\sf@size\usefont{U}{lasy}{m}{n}\symbol{41}}}}
\makeatother

\usepackage{tikz}
\usetikzlibrary{calc,angles,quotes,arrows.meta} % for pic, angle labels and arrow size
\tikzset{>=latex} % for LaTeX arrow head
\usepackage[outline]{contour} % glow around text
\contourlength{1.2pt}

\colorlet{wall}{blue!30!black}
\colorlet{xcol}{blue!70!black}
\colorlet{vcol}{green!60!black}
\colorlet{myred}{red!65!black}
\colorlet{mygreen}{green!60!black}
\colorlet{mydarkred}{red!50!black}
\tikzstyle{mysmallarr}=[-{Latex[length=3,width=2]}]

\colorlet{metalcol}{blue!25!black!30!white}
\tikzstyle{vvec}=[->,vcol,very thick,line cap=round]
\tikzstyle{force}=[->,myred,thick,line cap=round]
\tikzstyle{Fproj}=[->,myred!50,thick,line cap=round]
\tikzstyle{rope}=[brown!30!black,double=brown!90!black,double distance=1.2,line width=0.4,line cap=round]
\tikzstyle{dark rope}=[brown!10!black,double=brown!60!black,double distance=1.2,line width=0.4]
\tikzstyle{metal}=[draw=metalcol!10!black,rounded corners=0.1,
  top color=metalcol,bottom color=metalcol!80!black,shading angle=10]
\tikzstyle{wood}=[draw=brown!80!black,rounded corners=0.1,
  top color=brown!80,bottom color=brown!80!black!80,shading angle=10]
\tikzstyle{width}=[<->] %{Latex[length=4,width=3]}-{Latex[length=4,width=3]}]
\newcommand{\vbF}{\vb{F}}
\def\wave#1#2{
  ({(#1-0.36)*\xmax},0) to[out=0,in=180,looseness=1] (#1*\xmax,#2*\A)
  to[out=0,in=180,looseness=1]++ (0.36*\xmax,#2*-\A)
}

\newcommand\rightAngle[4]{
  \pgfmathanglebetweenpoints{\pgfpointanchor{#2}{center}}{\pgfpointanchor{#3}{center}}
  \coordinate (tmpRA) at ($(#2)+(\pgfmathresult+45:#4)$);
  \draw[white,line width=0.6] ($(#2)!(tmpRA)!(#1)$) -- (tmpRA) -- ($(#2)!(tmpRA)!(#3)$);
  \draw[mydarkred] ($(#2)!(tmpRA)!(#1)$) -- (tmpRA) -- ($(#2)!(tmpRA)!(#3)$);
}
\newcommand\lineend[2]{
  \def\w{0.1} \def\c{30}
  \draw[mygreen] (#1)++(#2:\w) to[out=#2-180-\c,in=#2+\c] (#1)
                               to[out=#2+\c-180,in=#2-\c]++ (#2-180:\w);
}


\title{Formulaire d'Onde S4}
\author{Raphaël Jamann}
\date{}

\begin{document}
    \maketitle

    \begin{multicols}{2}[
    \section*{Démonstration de l'équation de D'Alembert}
    ]
        
        \subsection*{Hypothèses}

            On considère une corde vibrante de masse linéique $\mu$, soumise à une force de tension $\vec F(x, t)$ dont la norme est $T(x, t)$.

            On suppose :
            \begin{itemize}
                \item \textbf{Hypothèse 1 (norme des tensions Cte) :} $\norm{\vec F_g(x,t)} = \norm{\vec F_d(x+dx,t)} = T_0$. \\
                        On retrouve cette hypothèse en projetant le PFD sur $(Ox)$.
                \item \textbf{Hypothèse 2 (petits angles) :} $\sin(\alpha) \approx \alpha$ et $\tan(\alpha) \approx \alpha \approx \frac{\partial y}{\partial x}$.
                \item Le mouvement est transversal : $y(x,t)$ désigne le déplacement vertical et $\alpha(x,t)$ l'angle avec l'horizontale.
            \end{itemize}

        \subsection*{Étude d'un élément de corde}

            PFD appliqué au système \{tronçon\} de masse négligeable dans un référentiel terrestre supposé galiléen :
            $$ \vec F_g(x,t) + \vec F_d(x + dx,t) = \mu dx \vec a $$

        \subsection*{Projection du PFD selon $Oy$}

            En utilisant l'approximation des petits angles ($\sin\alpha \approx \alpha$) :
            $$ T_0(\alpha(x+dx, t)-\alpha(x, t)) = \mu dx \frac{\partial^2 y(x, t)}{\partial t^2} $$

        \subsection*{Développement de Taylor et simplification}

            Par développement de Taylor à l'ordre 1 et simplification par $dx$ :
            $$ T_0 \frac{\partial \alpha(x,t)}{\partial x} dx = \mu dx \frac{\partial^2 y(x,t)}{\partial t^2} $$

            Or, comme $\tan(\alpha) \approx \alpha$, on a $\alpha(x,t) = \frac{\partial y(x,t)}{\partial x}$. D'où :
            $$ T_0 \frac{\partial^2 y(x,t)}{\partial x^2} = \mu \frac{\partial^2 y(x,t)}{\partial t^2} $$

        \subsection*{Équation de D'Alembert}

            On obtient la forme classique de l'équation de propagation de l'onde (ou équation de d'Alembert) à condition de poser $c = \sqrt{\frac{T_0}{\mu}}$ :
            $$
            \boxed{
            \frac{\partial^2 y(x,t)}{\partial x^2} = \frac{1}{c^2} \frac{\partial^2 y(x,t)}{\partial t^2}
            }
            $$
    \end{multicols}

    % TENSION - zoom
    \bigskip
    \begin{center}
        \begin{tikzpicture}[scale=2]
                \def\v{0.25*\xmax}
                \def\om{360/\xmax} % omega (degrees)
            \def\A{2.2}
            \def\xmax{10}
            \def\xa{0.87*\xmax}
            \def\xb{0.95*\xmax}
            \def\F{1.2} % tension
            \coordinate (L) at (\xa,{\A*cos(\om*\xa)});
            \coordinate (R) at (\xb,{\A*cos(\om*\xb)});
            \draw[dark rope,double distance=3.2,line width=1.0,samples=50,smooth,variable=\x,domain=\xa:\xb]
                plot(\x,{\A*cos(\om*\x)});
            \node[above left] at ({(\xb+\xa)/2},{\A*cos(\om*(\xb+\xa)/2)}) {$\dd{m}$};
            \draw[force,very thick] (\xa,{\A*cos(\om*\xa)}) --++ ({-atan(\A*2*pi/\xmax*sin(\om*\xa))-180}:\F)
                coordinate (T1) node[above=0,left=-2] {$\vec F(x,t)$};
            \draw[force,very thick] (\xb,{\A*cos(\om*\xb)}) --++ ({-atan(\A*2*pi/\xmax*sin(\om*\xb))}:\F)
                coordinate (T2) node[above=0,right=-2] {$\vec F(x+dx,t)$};
            \draw[dashed]
                (L) --++ (-\F,0) coordinate (X1)
                (R) --++ ( \F,0) coordinate (X2);
            \draw[width] (L)++({-atan(\A*2*pi/\xmax*sin(\om*\xa))-101}:0.1) --++ (\xb-\xa,0)
                node[midway,below] {$\dd{x}$};
            \draw[width] (R)++({-atan(\A*2*pi/\xmax*sin(\om*\xb))-79}:0.1) --++ (0,{\A*cos(\om*\xa)-\A*cos(\om*\xb)})
                node[midway,right] {$\dd{y}$};
            \draw pic[-,"{$\alpha(x,t)$}"{scale=0.7, above=-5},draw=black,angle radius=15,angle eccentricity=2]
                {angle=X1--L--T1};
            \draw pic[-,"{\contour{white}{$\alpha(x+dx,t)$}}"{scale=0.7,below=-3},draw=black,angle radius=15,angle eccentricity=2.8]
                {angle=X2--R--T2};
        \end{tikzpicture}
    \end{center}


    \clearpage

    \section{Solutions de l'équation de D'Alembert}

        \subsection{Forme générale (Ondes progressives)}
            La solution générale de l'équation de D'Alembert est la superposition de deux signaux se propageant sans déformation à la célérité $c$ :
            $$ \boxed{y(x,t) = f(x - ct) + g(x + ct)} $$

        \subsection{Ondes Planes Progressives Harmoniques (OPPH)}
            Une solution particulière fondamentale est l'onde sinusoïdale (ou harmonique). Pour une onde se propageant vers les $x$ croissants :
            $$ \boxed{y(x,t) = A \cos(\omega t - kx + \phi)} $$
            Les grandeurs caractéristiques sont :
            \begin{itemize}
                \item \textbf{Amplitude} $A$ : valeur maximale du déplacement.
                \item \textbf{Pulsation :} $\omega = \frac{2\pi}{T} = 2\pi f$
                \item \textbf{Nombre d'onde :} $k = \frac{2\pi}{\lambda} = \frac{\omega}{c}$
                \item \textbf{Phase} $(\omega t - kx + \phi)$ : la phase instantanée, avec $\phi$ la phase à l'origine.
            \end{itemize}
            \textbf{Relations clés :}
            $$ \boxed{\lambda = cT = \frac{c}{f}} \quad ; \quad \boxed{\omega = ck} $$

        \subsection{Relation de dispersion}
            L'injection de la solution harmonique dans l'équation de D'Alembert impose une relation entre les caractéristiques temporelles et spatiales :
            $$ \boxed{\omega = c \cdot k} $$
           
    \section{Équations de Maxwell}
        Les équations de Maxwell décrivent le comportement des champs électrique $\va{E}$ et magnétique $\va{B}$ dans la matière ou le vide.

        \begin{itemize}
            \item \textbf{Maxwell-Gauss :} $\div \va{E} = \frac{\rho}{\varepsilon_0}$
            \item \textbf{Maxwell-Faraday :} $\rot \va{E} = -\pdv{\va{B}}{t}$
            \item \textbf{Maxwell-Thomson (Flux) :} $\div \va{B} = 0$
            \item \textbf{Maxwell-Ampère :} $\rot \va{B} = \mu_0 \va{j} + \mu_0 \varepsilon_0 \pdv{\va{E}}{t}\, , \qquad (\mu_0\varepsilon_0 = \frac{1}{c^2})$
        \end{itemize}

        \subsection{Propagation dans le vide}
            Dans le vide en l'absence de sources ($\rho=0$ et $\va{j}=\va{0}$), on applique l'opérateur rotationnel à l'équation de Maxwell-Faraday :
            $$ \rot(\rot \va{E}) = \rot(-\pdv{\va{B}}{t}) = -\pdv{(\rot \va{B})}{t} $$
            En utilisant l'identité vectorielle $\rot(\rot \va{E}) = \grad(\div \va{E}) - \Delta \va{E}$ :
            $$ \grad(\div \va{E}) - \Delta \va{E} = -\pdv{(\rot \va{B})}{t} $$
            Or, dans le vide, $\div \va{E} = 0$ (Maxwell-Gauss) et $\rot \va{B} = \mu_0 \varepsilon_0 \pdv{\va{E}}{t}$ (Maxwell-Ampère) :
            $$ \va{0} - \Delta \va{E} = -\mu_0 \varepsilon_0 \pdv[2]{\va{E}}{t} \implies \Delta \va{E} - \frac{1}{c^2} \pdv[2]{\va{E}}{t} = \va{0} $$
            Les champs électrique et magnétique vérifient donc l'équation de D'Alembert :
            $$ \boxed{ \Delta \va{E} = \frac{1}{c^2} \pdv[2]{\va{E}}{t} } \quad \text{et} \quad \boxed{ \Delta \va{B} = \frac{1}{c^2} \pdv[2]{\va{B}}{t} } $$
            où $c = \frac{1}{\sqrt{\mu_0 \varepsilon_0}} \approx 3 \cdot 10^8 \text{ m.s}^{-1}$ est la célérité de la lumière.

        \subsection{Structure de l'onde plane harmonique}
            Pour une OPPH de vecteur d'onde $\va{k}$ et de pulsation $\omega$ se propageant dans le vide, les équations de Maxwell imposent :
            \begin{itemize}
                \item \textbf{Transversalité :} $\va{k} \vdot \va{E} = 0$ et $\va{k} \vdot \va{B} = 0$. Les champs sont perpendiculaires à la direction de propagation.
                \item \textbf{Relation de structure :} $\va{B} = \frac{\va{k} \wedge \va{E}}{\omega}$
                \item \textbf{Propriétés :} Les champs $\va{E}$ et $\va{B}$ sont en phase, perpendiculaires entre eux, et vérifient $\norm{\va{E}} = c \norm{\va{B}}$. Le trièdre $(\va{k}, \va{E}, \va{B})$ est orthogonal direct.
            \end{itemize}

        \subsection{Aspect énergétique et Intensité}
            \begin{itemize}
                \item \textbf{Vecteur de Poynting :} $\va{\Pi} = \frac{\va{E} \wedge \va{B}}{\mu_0}$ (W.m$^{-2}$). Il donne la direction et le flux de puissance.
                \item \textbf{Densité volumique d'énergie :} $u_{em} = \frac{\varepsilon_0 E^2}{2} + \frac{B^2}{2\mu_0}$. Dans le vide, il y a équipartition : $u_e = u_m$.
                \item \textbf{Intensité (moyenne temporelle) :} Pour une OPPH :
                $$ \langle \Pi \rangle = I = \frac{E_0^2}{2\mu_0 c} = \frac{1}{2}\varepsilon_0 c E_0^2 $$
                \item \textbf{Impédance du vide :} $Z_0 = \sqrt{\frac{\mu_0}{\varepsilon_0}} = \frac{E}{H} \approx 377 \, \Omega$ (avec $\va{B} = \mu_0 \va{H}$).
            \end{itemize}

        \subsection{Relations de passage}
            À l'interface entre deux milieux (1) et (2), de normale unitaire $\va{n}_{1 \to 2}$, les champs subissent des discontinuités si des charges ou courants de surface sont présents :
            \begin{itemize}
                \item \textbf{Composantes normales :} 
                $$ E_{2n} - E_{1n} = \frac{\sigma}{\varepsilon_0} \quad \text{et} \quad B_{2n} - B_{1n} = 0 $$
                \item \textbf{Composantes tangentielles :}
                $$ \va{E}_{2t} - \va{E}_{1t} = \va{0} \quad \text{et} \quad \va{B}_{2t} - \va{B}_{1t} = \mu_0 (\va{j}_s \wedge \va{n}_{1 \to 2}) $$
            \end{itemize}
            Où $\sigma$ est la densité surfacique de charge et $\va{j}_s$ la densité surfacique de courant.
    \clearpage

    \section{Interférences}

        \subsection*{différence de chemin optique et différence de marche}

            La \textbf{différence de chemin optique} entre deux ondes arrivant en un point $M$ est définie par :
            $$ \delta = n_1 S_1 + n_2 S_2 + \dots + n_k S_k $$
            où $n_i$ est l'indice du milieu traversé et $S_i$ la distance parcourue dans ce milieu.

            La \textbf{différence de marche} est la différence de chemin optique exprimée en nombre d'ondes :
            $$ \Delta = \frac{\delta}{\lambda} $$


        % TWO SLIT PATH DIFFERENCE
        \begin{center}
        \begin{tikzpicture}[scale=1.5]
        
            \def\L{5.9}       % distance between walls
            \def\H{3.0}       % total wall height
            \def\f{0.9}       % fractional height of projection point
            \def\ang{atan((\f*\H+\a)/\L/2)} % theta
            \def\t{0.15}      % wall thickness
            \def\a{1.5}       % slit distance
            \def\d{0.20}      % slit size
            \coordinate (T) at (0,\a/2);
            \coordinate (B) at (0,-\a/2);
            \coordinate (L) at (0,0);
            \coordinate (R) at (\L,0);
            \coordinate (P) at (\L,\f*\H/2);
            \coordinate (M) at ($(B)!(T)!(P)$);
            \fill[mygreen!80!black] (T) circle (0.05) node[above=5 ,left=5] {$S_1$};
            \fill[mygreen!80!black] (B) circle (0.05) node[below=5 ,left=5] {$S_2$};
            
            % LINES
            \draw[mygreen,thick] (T) -- (P) node[midway,above=-1] {$r_1$};
            \draw[mygreen,thick] (B) -- (P) node[midway,below=3,right=6] {$r_2$}; %right=6,below right=-4
            \draw[dashed] (L) -- (P);
            \draw[dashed,black!60] (L) -- (R);
            
            % ANGLES
            \coordinate (Bshift) at ($(B)+({\ang-90}:0.1)$);
            \coordinate (Mshift) at ($(M)+({\ang-90}:0.1)$);
            \draw pic[mysmallarr,"$\theta$",mydarkred,draw=mydarkred,angle radius=38,angle eccentricity=1.14]
                {angle = R--L--P};
            
            % MEASURES
            \draw[<->,black] (0,-0.47*\H) --++ (\L,0) node[midway,fill=white,inner sep=1] {$D$};
            \draw[<->,black] (-2.1*\t,-\a/2) --++ (0,\a) node[midway,fill=white,inner sep=1] {$a$};
            \draw[<->,black] ([shift={(2.1*\t,0)}]P) -- ([shift={(2.1*\t,0)}]R)
                node[midway,fill=white,inner sep=1]{$x$};
            
            % WALL
            \fill[wall]
                (0,\a/2-\d/2) rectangle (-\t,-\a/2+\d/2)
                (0,\a/2+\d/2) rectangle (-\t,\H/2)
                (0,-\a/2-\d/2) rectangle (-\t,-\H/2)
                (\L,-\H/2) rectangle (\L+\t,\H/2);
            \fill[mygreen!80!black] (P) circle (0.3*\t) node[right=1,above left=-2] {P};
            
        \end{tikzpicture}
        \end{center}

    \subsection*{Démonstration de la formule de la différence de marche $\Delta = \frac{a x}{D}$}

        On fait deux hypothèses :
        \begin{itemize}
            \item \textbf{Hypothèse 1 :} $D \gg a$ (distance entre les fentes et l'écran est grande devant la distance entre les fentes).
            \item \textbf{Hypothèse 2 :} $D \gg x$ (distance entre les fentes et l'écran est grande devant la position du point d'observation sur l'écran, petits angles).
        \end{itemize}

        En développant $r_2^2 - r_1^2$ à l'aide du théorème de Pythagore et en utilisant les hypothèses d'approximation, on trouve :
        \begin{align*}
            r_2^2 - r_1^2 &= (D^2 + (x + \frac{a}{2})^2) - (D^2 + (x - \frac{a}{2})^2) \\
            &= (x + \frac{a}{2})^2 - (x - \frac{a}{2})^2 \\
            (r_2 - r_1) (r_2 + r_1) &= 4 x \cdot \frac{a}{2} = 2 a x, \quad \text{or } r_2 + r_1 \approx 2D \\
            \implies r_2 - r_1 &\approx \frac{2 a x}{2 D} = \frac{a x}{D}
        \end{align*}

    \subsection{Autre démonstration}
            
    \begin{minipage}[c]{0.45\textwidth}
        $$ \delta= n_0 a \sin{(\theta)} \approx n_0 a \theta $$
        Or 
        $$ \tan{(\theta)} \approx \theta = \frac{x}{D} $$
        Donc 
        $$ \delta = \frac{n_0 a x}{D} = n_0(r_2- r_1)$$
    \end{minipage}\hfill
    \begin{minipage}[c]{0.5\textwidth}
        \centering
        % TWO SLIT PATH DIFFERENCE close up
        \begin{tikzpicture}[scale=1.5]
            
            \def\L{5.5}       % distance between walls
            \def\l{3.8}       % distance between walls
            \def\H{3.5}       % total wall height
            \def\f{0.9}       % fractional height of projection point
            \def\t{0.15}      % wall thickness
            \def\a{1.6}       % slit distance
            \def\d{0.20}      % slit size
            \def\ang{27}      % angle
            \coordinate (T) at (0,+\a/2);
            \coordinate (B) at (0,-\a/2);
            \coordinate (L) at (0,0);
            \coordinate (R) at (\L,0);
            \coordinate (I) at ({\a/2/tan(\ang)},0);
            
            % LINES
            \draw[mygreen,thick] (T) --++ (\ang:.8*\l) coordinate (PT) node[midway,below=1,above left=-2] {$r_1$};
            \draw[mygreen,thick] (B) --++ (\ang:\l) coordinate (PB) node[midway,left=2,below right=-1] {$r_2$};
            \draw[mydarkred,dashed] (T) -- ($(B)!(T)!(PB)$) coordinate (M);
            \draw[black!60,dashed] (L) --++ (\ang:.9*\l) coordinate (PR);
            %\draw[black!60,dashed] (T) --++ (.20*\L,0) coordinate (TR);
            \draw[black!60,dashed] (L) --++ (.6*\L,0) coordinate (LR);
            %\draw[black!60,dashed] (B) --++ (.20*\L,0) coordinate (BR);
            
            % LINE END
            \lineend{PT}{\ang+70}
            \lineend{PB}{\ang+70}
            
            % ANGLES
            \draw pic[mysmallarr,"\contour{white}{$\theta$}",mydarkred,draw=mydarkred,angle radius=18.8,angle eccentricity=1.38 ]
                {angle = B--T--M};
            %\draw pic[->,"$\theta$",mydarkred,draw=mydarkred,angle radius=16,angle eccentricity=1.41]
            %  {angle = TR--T--PT};
            %\draw pic[->,"$\theta$",mydarkred,draw=mydarkred,angle radius=16,angle eccentricity=1.41]
            %  {angle = BR--B--PB};
            \draw pic[mysmallarr,"$\theta$",mydarkred,draw=mydarkred,angle radius=22,angle eccentricity=1.25]
                {angle = LR--L--PR};
            \draw pic[mysmallarr,"$\theta$",mydarkred,draw=mydarkred,angle radius=22,angle eccentricity=1.30]
                {angle = LR--I--PB};
            %\rightAngle{T}{M}{PB}{0.3}
            % angle droit entre t M et PB :
            \draw[mydarkred] ($(M)!0.15!(T)$) -- ($(M)!0.15!(T)+({\ang}:0.2)$)--($(M)!0.07!(PB)$);

            
            % MEASURES
            \draw[<->,black] (-2.2*\t,-\a/2) --++ (0,\a) node[midway,fill=white,inner sep=1.5] {$a$};
            \draw[<->,black] ([shift={({\ang-90}:0.1)}]B) -- ([shift={({\ang-90}:0.1)}]M)
                node[midway,above=1,below right=-3]{$a\sin\theta$};
            
            % WALL
            \fill[wall]
                (0,\a/2-\d/2) rectangle (-\t,-\a/2+\d/2)
                (0,\a/2+\d/2) rectangle (-\t,\H/2)
                (0,-\a/2-\d/2) rectangle (-\t,-\H/2);
            
        \end{tikzpicture}
    \end{minipage}

    
    \subsection*{Calcul intensité}
        $$ I_{\text{tot}} = A^2 = I_1 + I_2 + 2\sqrt{I_1 I_2} \cos(\Delta \Phi) $$
        \begin{itemize}
            \item $I_1 = \left(\dfrac{E_{01}}{r_m}\right)^2$ et $I_2 = \left(\dfrac{E_{02}}{r_m}\right)^2$ représentent l'intensité reçue en $P$ lorsqu'une seule source émet.
            \item Si $E_{01} = E_{02} = E_0$, alors $I_{\text{tot}} = 2I_0^2 (1 + \cos(\Delta \Phi)) $.
            \item $\Delta \Phi = k\,(r_2 - r_1) + (\varphi_1 - \varphi_2) = \dfrac{2\pi}{\lambda}\left((r_2 - r_1) + (\varphi_1 - \varphi_2)\right)$ représente le déphasage entre les deux ondes reçues en $P$.
            \item Ce déphasage est la somme de deux termes : le premier terme rend compte de la différence $r_2 - r_1$ des trajets géométriques des deux ondes pour aller jusqu'en $P$ ; le deuxième représente la différence de phase $(\varphi_1 - \varphi_2)$ entre les ondes au moment de leur émission. Pour des sources cohérentes, ce terme ne dépend pas du temps.
            \item $2\sqrt{I_1 I_2} \cos(\Delta \Phi)$ est le terme interférentiel qui rend compte de la variation de l'intensité en fonction du point de mesure.
        \end{itemize}
        Pour calculer $A^2$, on utilise $A^2 = \underline{E}_{\text{tot}} \times \overline{\underline{E}_{\text{tot}}}$
    
    \subsection{Ordre d'Interférence}
        
        \begin{itemize}
            \item Les interférences sont constructives (maximum d'intensité) lorsque $\Delta \Phi = 2 p \pi$ (avec $p \in \mathbb{Z}$), soit lorsque $\delta = p\lambda_v$ (différence de marche entière).
            \item Les interférences sont destructives (minimum d'intensité) lorsque $\Delta \Phi = (2 p + 1) \pi$ (avec $p \in \mathbb{Z}$), soit lorsque $\delta = (p + \frac{1}{2})\lambda_v$ (différence de marche demi-entière).
        \end{itemize}

        On calcul l'ordre d'interférence $p$ à partir de $\delta$ et de $\lambda_v$ :
        $$ p = \frac{\delta}{\lambda_v} $$        

\end{document}